\documentclass[11pt]{article}
\usepackage[margin=1in]{geometry}
\usepackage[T1]{fontenc}
\usepackage{lmodern}
\usepackage{amsmath,amssymb}
\usepackage{booktabs}
\usepackage[numbers,sort&compress]{natbib}
\usepackage[hidelinks]{hyperref}
\usepackage{graphicx}
\sloppy

\title{Counting the actions no log recorded:\\capture-recapture for agent trace completeness}
\author{Caio Theodoro}
\date{October 2026}

\begin{document}
\maketitle

\begin{abstract}
Agent audit trails are engineered for integrity, not completeness: nothing in
the stack estimates how many consequential actions are missing from the record.
When no single log is complete, that count is a population-size problem.
Capture-recapture over three logs written by different infrastructure --- the
agent's own trace, a gateway proxy log and an end-of-shift environment diff ---
estimates the actions none of them saw. I test this with a benign live agent in
two environments whose kernels keep a ledger no log reads, with log loss
simulated afterwards. Two pre-registered tests of a BIC-selected log-linear
estimator missed their coverage bar under the record linkage they registered,
matching records by effect: 0.893 in expectation against a 0.90 bar in
back-office, 0.401 in infra. With exact record matching, which no auditor has,
both are nominal (0.948, 0.946). Post hoc, linking the trace and the gateway by
the tool-call id they share restores 0.946 in back-office but only 0.859 in
infra. Two further limits are the estimator's own. Loss that concentrates on the
actions the environment diff cannot see makes it fail silently, and loss in
bursts does too, mostly because the agent works one kind of ticket at a time and
a lost block lands where the diff is blind. A dispersion correction trades width
for coverage there; it does not remove the bias. A full technical report, every
recorded episode and every artifact behind these numbers are released with it.
\end{abstract}

\section{Integrity is engineered; completeness is assumed}

Audit regimes that reach AI agents start from a log~\citep{euaiact}, and the
engineering around that log answers one question: whether a line was changed
after it was written~\citep{flightrecorder2026,vet2025}. None of it answers
whether the record is whole. Agents have been shown deleting their own trace
events and spoofing tool calls~\citep{qin2026tamper}. A deleted line in a hash
chain is detectable; a line never written is not.

The count of actions no log shows needs estimating only when no single record is
complete. If every action passes one channel that numbers its action entries
before anything is dropped, the gaps give the count exactly. That fails when an
action bypasses the channel, when logs belong to parties that do not share their
numbering, or when sampling happens first. Epidemiology has a method for the rest:
overlapping several incomplete registers estimates what all of them
missed~\citep{fienberg1972,bird2018mse}. I found no work applying it to agent
trace completeness; the nearest, over LLM judges counting label errors, produced
valid intervals too wide to beat a raw count~\citep{kansal2026labels}.

This paper reports what happened when I pre-registered that test. The method
recovered the count only under exact record matching. Matching records the way
an auditor can is where it broke first, and two loss patterns break it under any
matching. Section~\ref{sec:setup} gives the setup, \S\ref{sec:prereg} the two
pre-registered tests, \S\ref{sec:linkage} the linkage failure,
\S\ref{sec:loss} the loss the estimator does not expect, and \S\ref{sec:limits}
what none of it shows.

\section{Setup}
\label{sec:setup}

An \emph{action} is a tool call that changes environment state or sends data
out. The environment kernel appends one ledger entry per action, so the true
count $N$ is a program, not a label, and no witness reads it. Three witnesses
observe actions: W1, the harness's trace of the agent's tool calls; W2, a gateway
proxy log; W3, an end-of-shift diff of the environment's state plus its outbound
mail sink. Each action gets a capture history, and the cell for ``seen by none''
is what is estimated.

The estimator, \texttt{ll\_bic}, fits Poisson log-linear models to the seven
observed cells~\citep{fienberg1972}, selects the pairwise interactions by BIC
and reports a 95\% profile-likelihood interval. Two baselines bound it: the union
of all three logs, which never covers once anything is hidden, and the interval
$[n, \infty)$, which always covers and says nothing. Gap detection is the share of
replicates whose lower end exceeds the union: the estimator asserting the logs
are incomplete.

Records are matched across logs in one of two ways. Exact matching uses the
ledger id, which no auditor has. Effect-signature matching pairs records on
operation, resource and, for some operations, value, in order within a
signature.

A live agent (\texttt{deepseek-v4.1-flash}, benign) worked two environments: a
back-office queue of refunds, cancellations, customer updates and email, and an
infra/ops queue whose end-state diff misses transient work. Before any loss, the
trace and the gateway each held every action, so all loss is simulated
afterwards, per record, on the recorded episodes. In back-office the env diff
sees only financial operations, by configuration; in infra its blind spots come
from the domain. A full description, including the LLM auditor that turned out to
be a copy of the trace, is in the technical report.

\section{Two pre-registered tests}
\label{sec:prereg}

Both tests were frozen before their paid runs, on loss configurations held out of
the estimator choice, with effect-signature matching as the headline
(Table~\ref{tab:prereg}).

\begin{table}[h]
\centering
\resizebox{\linewidth}{!}{\begin{tabular}{llrrrrrr}
\toprule
pool & linkage & \texttt{ll\_bic} & bias & \texttt{ll\_independence} & W1:W2 alone & its coverage & its bias \\
\midrule
registered, 21 & oracle & 0.948 $\pm$ 0.004 & +0.31 & 0.951 & 0.010 & 0.867 & -0.71 \\
 & signature & 0.893 $\pm$ 0.006 & +6.20 & 0.948 & 0.358 & 0.858 & +12.71 \\
restored, 24 & oracle & 0.945 $\pm$ 0.004 & -0.31 & 0.951 & 0.008 & 0.783 & -3.14 \\
 & signature & 0.862 $\pm$ 0.006 & +8.97 & 0.949 & 0.517 & 0.847 & +13.41 \\
\bottomrule
\end{tabular}}
\par\smallskip{\footnotesize \texttt{fin\_mid}, 3000 replicates per cell (30 fresh loss seeds), both linkages on the same random numbers. ``W1:W2 alone'' is the share of replicates where BIC selected the model with only the trace--gateway interaction; ``its coverage'' and ``its bias'' are \texttt{ll\_bic}'s on those replicates. Post hoc. Source: \texttt{artifacts/robustness/precision.json}.}

\caption{Back-office, held-out per-record loss, by linkage. Post hoc precision on
the registered episodes.}
\label{tab:prereg}
\end{table}

In back-office the coverage hypothesis passed at its registered loss seed (0.940
against a 0.90 bar), with the gap detected in every replicate. At 30 fresh loss
seeds the same analysis covers 0.893 under signature matching and 0.948 under
exact matching; restoring three episodes the run lost gives 0.862 and 0.945. In
infra the transfer hypothesis failed outright: 0.300 at the registered seed, 0.401
over 20 seeds, against 0.946 exact (Table~\ref{tab:transfer}). In both
environments the union never covers, the estimator asserts the logs are
incomplete, and when the gateway and the trace fail together the independence
model collapses while the selected one holds.

\begin{table}[h]
\centering
\resizebox{\linewidth}{!}{\begin{tabular}{llrrrrr}
\toprule
config & linkage & \texttt{ll\_bic} & per seed & bias & width & \texttt{ll\_independence} \\
\midrule
\texttt{transfer\_mid} & oracle & 0.946 $\pm$ 0.005 & 0.89--0.99 & -0.18 & 59.8 & 0.948 \\
 & signature & 0.401 $\pm$ 0.011 & 0.28--0.48 & -32.62 & 55.8 & 0.346 \\
\texttt{transfer\_mid\_rho100} & oracle & 0.945 $\pm$ 0.005 & 0.91--0.98 & +2.20 & 118.5 & 0.000 \\
 & signature & 0.799 $\pm$ 0.009 & 0.73--0.87 & -32.16 & 109.6 & 0.000 \\
\bottomrule
\end{tabular}}
\par\smallskip{\footnotesize Infra/ops environment, paid run: 23 episodes, N=1036. 2000 replicates per cell (20 fresh loss seeds), both linkages on the same random numbers. The registered T1 cell (one seed) is 0.300. Source: \texttt{artifacts/robustness/transfer-precision.json}.}

\caption{Infra, by linkage. Post hoc precision on the registered episodes.}
\label{tab:transfer}
\end{table}

\section{Record linkage is where it breaks first}
\label{sec:linkage}

The two environments fail differently. In back-office, 0.223 of actions share an
effect signature with another action in the same episode, almost all of them
emails to the same customer. In-order pairing lines up records of different
actions as one action seen by both the trace and the gateway, and BIC fits that
fake dependence in 0.358 of replicates against 0.010 under exact matching. In
infra, repeated identical actions merge away and the env diff is mis-attached.

The same mechanism produces false gaps. When the gateway and the trace lose
mutually exclusive actions, so that nothing is hidden, signature matching claims
a gap at every drop rate and in every pool tested once the env diff is partial;
exact matching claims none.

The signature was a poor design choice. It left out email and file bodies, which
every log carries as a digest, and the trace and the gateway share the provider's
tool-call id. Table~\ref{tab:arms} compares the linkages an auditor can compute.
Linking trace and gateway by call id restores 0.946 in back-office, where the env
diff sees every email and is easy to attach, and removes every false gap measured,
but that is close to true by construction. In infra it reaches only 0.859: the env
diff carries no id, and attaching its rows to repeated identical actions needs a
per-row id or timestamp the diff does not have. The call id is safe only while
infrastructure writes it; an agent writing its own trace could forge it, and that
was not tested.

\begin{table}[h]
\centering
\resizebox{\linewidth}{!}{\begin{tabular}{llrrrrr}
\toprule
setting & config / score & signature & digest & call id & last writer & exact \\
\midrule
back-office, financial view & \texttt{fin\_mid} & 0.875 & 0.942 & 0.946 & 0.946 & 0.946 \\
same, restored pool & \texttt{fin\_mid} & 0.852 & 0.944 & 0.945 & 0.945 & 0.945 \\
back-office, full view & \texttt{full\_mid} & 0.921 & 0.905 & 0.924 & 0.922 & 0.940 \\
same, exploratory pool, sampled & \texttt{full\_sampled} & 0.917 & 0.930 & 0.921 & 0.931 & 0.941 \\
infra & \texttt{transfer\_mid} & 0.406 & 0.379 & 0.859 & 0.888 & 0.944 \\
infra, coupled & \texttt{transfer\_mid\_rho100} & 0.803 & 0.803 & 0.803 & 0.839 & 0.944 \\
back-office, nothing hidden & false gap & 0.00--1.00 & 0.00--0.00 & 0.00--0.00 & 0.00--0.00 & 0.00--0.00 \\
 & coverage & 0.00--0.99 & 0.29--0.81 & 1.00--1.00 & 1.00--1.00 & 1.00--1.00 \\
infra, nothing hidden & false gap & 0.18--1.00 & 0.00--0.69 & 0.00--0.00 & 0.00--0.00 & 0.00--0.00 \\
 & coverage & 0.10--0.59 & 0.00--0.14 & 1.00--1.00 & 1.00--1.00 & 1.00--1.00 \\
\bottomrule
\end{tabular}}
\par\smallskip{\footnotesize Coverage of \texttt{ll\_bic} under five ways of building the capture table. ``digest'' keeps the content digest the logs already store; ``call id'' links trace and gateway by the provider's tool-call id and attaches the env diff by signature; ``last writer'' attaches the env diff's end-state row for an overwritten field to the latest logged call on it; ``exact'' is the ledger id, which no auditor has. Held-out cells: 2000 replicates, on loss seeds different from Table~\ref{tab:precision}. Nothing-hidden rows: ranges over exclusive-loss rates 0.2--0.5 (both env views in back-office), 1000 replicates each. Post hoc. Source: \texttt{artifacts/robustness/linkage-arms.json}.}

\caption{Linkages an auditor can compute against exact matching. The last-writer
arm is withdrawn: it produces false gaps. Post hoc.}
\label{tab:arms}
\end{table}

\section{Loss the estimator does not expect}
\label{sec:loss}

\paragraph{Loss that depends on the action.}
Raising both logs' drop rate on the operations the env diff cannot see breaks
\texttt{ll\_bic} silently, under either linkage: 0.647 at a tenth more loss and
0.187 at a fifth more, on intervals of ordinary width that still assert the gap
(Table~\ref{tab:dose}). Stratifying on the operation (\texttt{strat\_bic})
covers where it answers, refusing more often as the strata thin; at higher doses
nearly all its coverage is refusal.

\begin{table}[h]
\centering
\resizebox{\linewidth}{!}{\begin{tabular}{rrrrrrrrr}
\toprule
drop $\times$ & hidden & \texttt{ll\_bic} oracle & \texttt{ll\_bic} & width & \texttt{strat\_bic} & refuses & finite cover & finite $n$ \\
\midrule
1$\times$ & 62.3 & 0.943 & 0.887 & 43.1 & 0.943 & 0.07 & 0.954 & 280 \\
1.1$\times$ & 74.7 & 0.647 & 0.657 & 49.9 & 0.933 & 0.14 & 0.953 & 258 \\
1.2$\times$ & 90.7 & 0.187 & 0.197 & 69.0 & 0.957 & 0.34 & 0.970 & 197 \\
1.3$\times$ & 104.8 & 0.303 & 0.280 & 122.3 & 0.960 & 0.52 & 0.979 & 143 \\
1.4$\times$ & 122.5 & 0.517 & 0.520 & 196.8 & 0.967 & 0.71 & 0.953 & 86 \\
1.6$\times$ & 158.1 & 0.857 & 0.807 & 519.5 & 0.917 & 0.97 & 0.875 & 8 \\
1.8$\times$ & 201.8 & 0.907 & 0.920 & 2161.1 & 0.983 & 1.00 & -- & 0 \\
\bottomrule
\end{tabular}}
\par\smallskip{\footnotesize Both logs' drop rate multiplied for \texttt{update\_customer} and \texttt{write\_file}, the operations the financial env view cannot see; \texttt{fin\_mid} base, 300 replicates per dose. Signature linkage except the third column. ``finite cover'' is \texttt{strat\_bic}'s coverage among its bounded intervals. Exploratory. Source: \texttt{artifacts/robustness/mnar-dose.json}.}

\caption{Loss concentrated on the operations the env diff cannot see.
Exploratory.}
\label{tab:dose}
\end{table}

\paragraph{Loss in bursts.}
Dropping blocks of 12 consecutive actions per log, under exact matching, takes
\texttt{ll\_bic} to 0.588 in back-office and 0.772 in infra
(Table~\ref{tab:controls}). Most of that is not dispersion. The agent works one
kind of ticket at a time, so the actions the financial view cannot see sit in one
stretch of each episode, and a lost block that lands there hides actions no
witness then sees. Keeping the order but permuting which actions the diff sees
restores 0.875 to 0.940; on the long-episode pools \texttt{ll\_bic} had been
overcounting by 29.6 and 58.8.

\begin{table}[h]
\centering
\resizebox{\linewidth}{!}{\begin{tabular}{lrrrr}
\toprule
bursts, $w = 12$ & actual & diff's view permuted & order shuffled & \texttt{ll\_bic\_hom}, actual \\
\midrule
confirmatory pool & 0.588 (+4.0) & 0.875 (+1.1) & 0.914 (+1.0) & 0.966 \\
long-60 & 0.352 (+29.6) & 0.881 (+0.9) & 0.907 (+1.0) & 0.877 \\
long-120 & 0.328 (+58.8) & 0.940 (-0.3) & 0.939 (+0.3) & 0.846 \\
infra & 0.772 (+1.7) & 0.791 (+0.4) & 0.912 (+0.2) & 0.994 \\
\bottomrule
\end{tabular}}
\par\smallskip{\footnotesize Coverage (bias) of \texttt{ll\_bic} under blocks of 12 lost actions, with the episode as recorded; with which actions the env diff sees permuted across the episode, order kept; and with the order of actions shuffled. Infra has no configured view, so permuting it changes nothing there. Exact matching, 800 replicates per cell. Post hoc. Source: \texttt{artifacts/robustness/final-controls.json}.}

\caption{What makes block loss fail. Post hoc.}
\label{tab:controls}
\end{table}

A quasi-likelihood interval~\citep{wedderburn1974quasi,mccullagh1989glm} with the
dispersion taken from variation between episodes (\texttt{ll\_bic\_hom}) keeps
\texttt{ll\_bic}'s point and only widens the interval, so it never covers less. It
reaches 0.952 under back-office bursts and 0.994 in infra, at two and a half to
four times the width, but on the long-episode pools it covers 0.877 and 0.846 by
widening around a biased point. It is a trade, to be reported beside the plain
interval, not a repair. Whole-session outages and session quality shared by all
three logs, covered in the technical report, defeat both intervals unless the env
diff stays up.

\section{What this does not show}
\label{sec:limits}

\begin{itemize}
\item \textbf{No adversary.} Every live agent was benign; the loss was simulated
on complete logs. Nothing here is evidence about an agent that hides on purpose.
\item \textbf{The environments were mine.} The infra run reproduced a free scripted
run, so it tests the environment and the linkage, not the agent, and its verdict
depends on a collision density its ticket mix sets.
\item \textbf{Post hoc fixes.} The linkage arms and the dispersion correction were
found after the pre-registered runs, on the same episodes.
\item \textbf{Cheaper checks not compared.} Seeded canary actions and
sequence-number gaps measure log loss directly and were not run head to head. In
the back-office environment every action passes the gateway, so a gateway that
numbered its action entries before sampling would make this method unnecessary
there.
\end{itemize}

\section{Reproducing this}

The repository holds both environments, the estimators, every recorded episode,
the decision log with a dated entry for every prediction and every result that
contradicted one, and every artifact a number here comes from; a test fails if a
figure in this paper has no artifact behind it. All live runs together cost
\$5.92. The capture histories and ground truth are released as a dataset.

\bibliographystyle{plainnat}
\bibliography{refs}

\end{document}
